\SourceLeaf{132}
To determine the logarithm of the difference, one will find $\log a - \log b$ in the second or third column according as this difference is smaller or larger than $\log 2$. In the first case, when one sets $\log a - \log b = Q$, the result will be $\log(a-b) = \log b - P = \log a - R$. In the second, when one sets $\log a - \log b = \log R$, the required logarithm will be
\[
= \log b + P = \log a - Q.
\]

It is curious that \textsc{Leonelli} did not at once notice this way of determining the logarithm of the difference when the difference of the logarithms is smaller than $\log 2$. In the book itself he gives a different, more complicated procedure for this case. Only when \textsc{Delambre}'s report pointed out the deficiency was he led to this readily apparent way of using the table. He explains the procedure at great length in his remarks on the report.

Although we cannot withhold our general approval of \textsc{Leonelli}'s idea of simplifying logarithmic calculations by such a table, and indeed consider it desirable that such a table should actually be produced, we can attach little value to the rest of what \textsc{Leonelli} says on the subject. His explanation of how to use the table is unnecessarily long for so elementary a matter. We find no investigation that goes even a little deeper into the law governing the progression of the tables and of the differences, though such an investigation would be essential to deciding how far the tables should extend. There are, however, a few offhand remarks, some expressed rather confusedly, which suggest that \textsc{Leonelli} either made no such investigation or did so quite incorrectly. This includes, e.\,g., the thoroughly false claim that, when $P$ increases by a fixed increment of 0,0005 and approaches the value 6, the first differences of $Q$ will have only one or two significant digits left if $Q$ is expressed to 14 decimal places. A simple calculation shows that these differences retain at least five significant digits up to that point. Still less do they vanish when $P$ exceeds 6. Only at $P = 10,7$ will they have decrea\SourceLeaf{133}%
